\section{Related work}

\paragraph{Linear relational embeddings.}
\citet{hernandez2024linearity} brought the idea of
modeling relations as linear maps between concept representations
\citep{DBLP:conf/icml/PaccanaroH00} to transformer
language models \citep{vaswani2017attention}. Rather than
training such maps, they estimate them from a trained model by
approximating the mapping from a subject's hidden state to the model's
output with an affine function derived from the model's Jacobian. The
approximation is faithful for some relations and poor for others.
\citet{merullo2025linear} trace part of this
variation to pretraining data: LREs are more faithful when subjects and
objects frequently co-occur in the corpus (we compare against such
frequency measures in Appendix \ref{sec:corpus-freq-baselines}). A more fundamental question is
whether linear structure reflects what a model has learned at all, since
many parametrizations induce the same next-token distribution.
\citet{marconato2024all} refine the notion of
relational linearity and show that, under suitable conditions, it holds
either in all models that induce the same next-token distribution or in
none. \citet{valer2026relational} build on this
formulation with an empirical probe based on the Kullback--Leibler
divergence and find that relational linearity varies across models and
layers and with how the relation is phrased. These works evaluate and
explain linearity at the level of the relation. We ask instead how much
of it can be predicted from how the model represents the subject, using
a readout estimated without the relation or any of its objects. We also
identify object cardinality as a second factor that shapes LRE
faithfulness (§\ref{sec:residuals}).

\paragraph{Subject representations in factual recall.} Mechanistic studies
converge on the view that much factual knowledge is attached to the
subject before the relation is processed. Feed-forward layers operate as
key--value memories whose keys respond to patterns in the input and
whose values induce distributions over the output vocabulary
\citep{geva2022transformer}, and causal tracing
localizes factual associations in middle-layer MLPs at the last subject
token, where targeted edits can change individual facts
\citep{meng2022locating}. Building on these findings,
factual recall has been decomposed into stages: MLP layers enrich the
subject representation with many of its attributes, the relation is
propagated to the final position, and attention heads extract the
requested attribute
\citep{geva-etal-2023-dissecting, nanda2024fact}.
\citet{nanda2024fact} further report that the lookup
performed by early MLPs depends largely on the tokens of the subject's
name rather than on the preceding context.
\citet{10.1162/TACL.a.746} show that knowledge
acquisition tracks how often subject and answer co-occur, though
frequency alone does not fully explain learning.

\paragraph{How facts are stored.} Theoretical work offers two accounts of
how transformers store facts.
\citet{nichani2025understanding} model attention and MLP
weights as associative memories and prove that a one-layer transformer
can store a number of facts that scales linearly with its parameter
count, with storage traded off between self-attention and MLP
parameters. \citet{DBLP:journals/corr/abs-2605-12426} propose a
geometric alternative: in a controlled one-layer setting, subject
embeddings encode linear superpositions of their associated attributes,
and a small MLP acts as a relation-conditioned selector through ReLU
gating rather than as a key--value lookup. The geometric account is
close in spirit to our hypothesis. If an entity's attributes are already
superposed in its representation, a relation only needs to select among
them, and selection should succeed to the extent that the attribute is
prominently encoded.

\paragraph{Lenses.} Lenses read out intermediate hidden states in terms of
the output vocabulary and differ mainly in the map they apply: the
model's own unembedding
\citep{nostalgebraist2020interpreting}, an affine translator
learned per layer \citep{belrose2023eliciting}, or
the model's input--output Jacobian averaged over a large corpus and
applied to raw activations, with the aim of identifying representations
the model is poised to verbalize
\citep{gurnee2026workspace}.
Our contrastive variant evaluates the Jacobian at reference entities of
the same type and applies it to differences between entity
representations. This subtracts what the reference entities share, so
that the top-ranked tokens describe an individual entity rather than its
category. As a predictor of LRE faithfulness, the non-contrastive
Jacobian lens performs comparably (Appendix \ref{sec:shared-baselines}).
%\citet{DBLP:conf/icml/PaccanaroH00} introduce linear
%relational embeddings (LREs)
%and \citet{hernandez2024linearity} show that LREs are part
%of the mechanism by which transformer models answer factual questions.
%
%
%\citet{merullo2025linear} show that the frequency
%of subject-object cooccurrence
%in the
%pretraining corpus is an important determinant of
%linearity. In our work, we put forward the hypothesis that
%relational linearity is primarily a property of how the model
%represents the subject, independent of the relation (or object). We also
%identify object cardinality as an additional important
%variable for linearity. Finally, we analyze limitations of
%LRE for understanding linearity and make suggestions for how
%future factual recall datasets should be constructed for
%relational linearity research.
%
%\citet{geva-etal-2023-dissecting} proposes subject
%enrichment as a key mechanism for factual recall. We can
%interpret our contrastive Jacobian lens as a probe of
%subject enrichment.
%Crucially, we connect in this paper two different lines of
%research on the representation of relations in transformer
%models: we hypothesize that the type of subject enrichment that the Jacobian
%lens diagnoses gives rise to the linearity of relations
%found by \citet{hernandez2024linearity}.
%
%DISCUSS \cite{10.1162/TACL.a.746}
%
%DISCUSS \cite{DBLP:journals/corr/abs-2605-12426}

